\section{Conclusion}

%In this work, we identified and defined problem drift as performance degradation during an ongoing LLM-based MAD.
%We conducted experiments on ten datasets across generative, knowledge, reasoning, and instruction-following tasks to quantify the relevance of problem drift.
In this work, we identified \textbf{problem drift}, a performance degradation in MAD.
We conducted experiments on ten datasets across ten different tasks to quantify problem drift and found that it occurs across all tasks.
Our experiments showed that problem drift affects 7\% (StrategyQA) to 89\% (ETPC) of debates, but especially when discussing low-complexity and generative problems.
We proposed \judge\ and \policy\ as first baselines to detect and mitigate problem drift.
\policy\ reduced occurrence of problem drift, mitigating 30.6\% of problem drift for weaker agents.
We characterized eight error types through a human evaluation with eight human participants.
These errors were grouped into two main categories: temporal (i.e., lack of progress, low-quality feedback, low-quality engagement) and local (i.e., lack of clarity, task non-compliance, knowledge gap, logical error, linguistic error).
We published the code, the \dataset\ dataset, and the human annotations.
% We further propose \policy, which introduces on-demand feedback about the issues leading to problem drift and mitigates 31\% of occurences.

%Future work could investigate how agent training, e.g., through reinforcement learning \citep{wang2024adaptingllmagentsuniversal}, could be used to avoid problem drift during test-time.
%One could investigate the usage of \judge\ as a reward model.
%Meanwhile, the identified error types help to understand in which situations current agents struggle to maintain high-quality discourse.
%Agent training could also involve dynamic agent orchestration \citep{liu2024dynamicllmpoweredagentnetwork} and consensus-based \citep{ChenSB24a} or tree-search \citep{HuMYD24, KohMFS24} decision-making.

% What does it mean that multiple agents debating drift away from the initial task goal?
% Is it a natural property that these systems (similar to humans) try to explore new ideas and solutions?
% Currently, our findings give us reason to believe that, in some cases, agents drift away to recover stronger later on in the discussion. 
% Still, they sometimes become prone to simple temporal and local mistakes.
% Future work could explore the role of agents exploring new ideas that are not directly task-related at first but help deep reasoning. 
% Exciting avenues also lie in comparing the inner differences (or similarities) between human and agent debates.

Future work could explore the role of agents in discussing new ideas that are not directly task-related but help to reason deeply.
Other interesting directions lie in comparing the inner dynamics of human and agent debates, such as differences in explorativeness, conciseness, and argument structure.
One can question whether problem drift is a natural property that multi-agent systems present when exploring new ideas and solutions, similar to how humans explore different reasoning paths in discussions.
Our findings suggest that, in some cases, agents drift away to recover later on in the discussion. 
Still, these agents regularly become prone to simple temporal and local mistakes.


% The key aim of this study is to provide others with this new lens of problem drift to study MAD.

% Future Work
% We concentrated on a voting decision protocol for this study to have a definitive solution at the end of each turn \citep{YangDKH24}.
% However, future work could investigate problem drift for more dynamic multi-agent interaction.
% This could also involve consensus-based \citep{ChenSB24a} or tree-search \citep{HuMYD24, KohMFS24} decision-making protocols.

% \section{Future Work}
% We showed that \judge\ can be a valid approach to mitigate problem drift at test time. 
% While our initial investigation was bound by the available computational resources, future work could evaluate mitigation strategies on a larger scale and diversified tasks.
% More mitigation methods could be tested to potentially minimize the required compute by guiding agents during their response already.

% We concentrated on a voting decision protocol for this study to have a definitive solution at the end of each turn \citep{YangDKH24}.
% However, future work could investigate problem drift for more dynamic multi-agent interaction.
% This could also involve consensus-based \citep{ChenSB24} or tree-search \citep{HuMYD24, KohMFS24} decision-making protocols.

% As agents prefer shorter responses in conversational settings, future work could investigate the influence of similar characteristics with other attributes.
% Possibly, agents could abuse such latent preferences to reach their goals and influence the discussion to their liking.
% This is especially relevant for applications that aim to provide fair decision-making or safe environments \citep{MukobiRRS23}.

\section*{Limitations}

% Multi-agent systems have several agents that repeatedly contribute, leading to extensive conversations.
The large amount of generated text in MAD can hinder qualitative human analysis of their discussions, as annotation is expensive.
In this human study, we investigated only the drifting turn and its immediate predecessor.
This decision provides a reasonable space to search for local and short-context temporal mistakes agents make during the debate. 
Still, it does not consider the full discussion history, which may include valuable information.
To keep our findings on error types comparable to the actual MAD, we also limit the context length of the agents to the current and previous turns during all experiments.

Our \judge\ requires forwarding solution pairs from consecutive turns, which can be computationally expensive. 
The combination of \judge\ and \policy\ yields an 8.5\% increase in computational cost (cf. \Cref{tab:compute} of \Cref{app:compute}).
We did not find this limiting for this project, but note that for production environments, fine-tuning a more cost-efficient classifier like BERT \citep{sun2020finetuneberttextclassification} on drifting and non-drifting examples could be a promising alternative.

%There is always the possibility that in a different multi-agent setup, results differ.

\section*{Acknowledgements}
This work was partially supported by the Lower Saxony Ministry of Science and Culture and the VW Foundation. 
We acknowledge EuroHPC Joint Undertaking for awarding us access to MeluXina at LuxProvide, Luxembourg.
This work was funded by the Deutsche Forschungsgemeinschaft (DFG, German Research Foundation) – 564661959.
Many thanks to Dominik Meier for the thoughtful discussions.